\section{Introduction}
\label{sec:introduction}

Hikita's conjecture relates two different kinds of algebra: the cohomology
of a symplectic resolution and the functions on a fixed subscheme of a
corresponding Coulomb branch.  The scheme structure is essential.  Even
when the underlying fixed-point set is small, its nilpotent functions can
encode positive-degree cohomology.  For quiver gauge theories, the
tautological classes on the Higgs side and the equivariant parameters on
the Coulomb side give a common coefficient ring.  The question is whether
the relations among those coefficients agree.

We prove this agreement for arbitrary finite quivers, under the precise
regularity hypothesis stated below.  The proof keeps the coefficient maps
and the full fixed schemes throughout.  Its two main comparisons have
different purposes: a calculation for the full topological torus
identifies the cohomology algebra, while a separate argument identifies
that torus-fixed scheme with the fixed scheme of the stability
cocharacter.

\subsection{The statement and its coefficient maps}

Let $Q=(Q_0,Q_1)$ be a finite oriented quiver.  Loops and multiple arrows are
allowed.  Fix complex vector spaces $V_i,W_i$ of dimensions $v_i,w_i$, and put
\begin{equation}\label{eq:gauge-data}
 G=\prod_{i\in Q_0}\GL(V_i),\qquad
 N=\bigoplus_{a:i\to j}\Hom(V_i,V_j)
       \oplus\bigoplus_{i\in Q_0}\Hom(V_i,W_i).
\end{equation}
An algebraic torus $F$ acts linearly on $N$ and commutes with $G$; write
$\mathfrak f=\Lie F$.  The induced action on $T^*N=N\oplus N^*$ is Hamiltonian.
For $(s,p)\in N\oplus N^*$, our cotangent moment map is characterized by
$\langle\mu(s,p),\xi\rangle=\langle p,\xi s\rangle$.
Choose a character $\nu$ of $G$ for which the $\nu$-semistable locus in
$\mu^{-1}(0)$ equals the stable locus and $G$ acts freely there.  We call such
a character \emph{regular}, and set
\begin{equation}\label{eq:higgs-definition}
 D=\mu^{-1}(0)^{\nu\text{-stable}},\qquad X=D/G.
\end{equation}
Thus $X$ is a smooth Nakajima variety.  We allow $D$ and $X$ to be empty.

Put $\OO=\C[[t]]$ and $\KK=\C((t))$.  The Braverman--Finkelberg--Nakajima
space of triples and its flavor-deformed Coulomb algebra are
\begin{align}
 \cR&=\{[g,s]\in G(\KK)\times_{G(\OO)}N(\OO):gs\in N(\OO)\},
       \label{eq:triples-definition}\\
 A_F&=H_*^{G(\OO)\times F}(\cR;\C),\qquad
 Y_{\mathfrak f}=\Spec A_F.\label{eq:coulomb-definition}
\end{align}
Here $(g,s)\sim(gh^{-1},hs)$, and the multiplication is BFN convolution.
We use the commutative algebra, without the loop-rotation deformation
parameter.  Its topological grading is by
$\pi_1(G)=\Z^{\{i:v_i>0\}}$, and defines a torus
$T_{\mathrm{top}}=\Hom(\pi_1(G),\C^\times)$.
The character $\nu$ gives a cocharacter of this torus.  Its fixed locus is
scheme-theoretic: if $(A_F)_m$ denotes its weight-$m$ subspace, then
\begin{equation}\label{eq:nu-quotient}
 \C[Y_{\mathfrak f}^{\nu}]
       =A_F/\langle(A_F)_m:m\ne0\rangle.
\end{equation}
No reduction of this quotient will be taken.

Let $T\subset G$ be the diagonal maximal torus, let $\mathcal W$ be its Weyl
group, and write
\begin{equation}\label{eq:coefficient-ring}
 R=H^*_{G\times F}(\mathrm{pt};\C)
   =\C[\Lie T\oplus\mathfrak f]^{\mathcal W},\qquad
 R_0=H_G^*(\mathrm{pt};\C).
\end{equation}
The Kirwan map $\phi_H:R\to H_F^*(X;\C)$ and the coefficient map
$\phi_C:R\to\C[Y_{\mathfrak f}^{\nu}]$ are the natural maps on the two
sides of the Hikita conjecture.  On the first side, a gauge Chern class
maps to the Chern class of its tautological bundle, and a flavor parameter
acts through $H_F^*(\mathrm{pt})$.  On the second side, $R$ acts by the
coefficient classes of the Coulomb algebra, followed by the fixed-scheme
quotient.  Each linear gauge or flavor parameter has degree two.

\begin{theorem}\label{thm:main}
For every datum \eqref{eq:gauge-data}, every commuting flavor torus $F$, and
every regular stability character $\nu$, both coefficient maps are
surjective and have the same kernel.  They therefore induce a canonical
isomorphism of graded
$R$-algebras
\[
 H_F^*(X;\C)\ \cong\ \C[Y_{\mathfrak f}^{\nu}].
\]
The gradings are ordinary cohomological degree and BFN homological degree,
respectively.  The assertion includes the empty case, whose ring is zero.
\end{theorem}

This is the equivariant cohomological Hikita conjecture for quiver gauge
theories, in the form of \cite[Conjecture~2.1, Equation~(2.3)]{DK25}.
It concerns the full coefficient-algebra quotient, with its nilpotents.
No finite-type, Dynkin, conical, or good/ugly hypothesis is imposed.

\subsection{History and the comparison problem}

A classical model is the description of type-$A$ Springer-fiber
cohomology through scheme-theoretic intersections of nilpotent orbit
closures with a Cartan subalgebra, developed by de Concini--Procesi and
Tanisaki \cite{DeConciniProcesi81,Tanisaki82}.  In the setting of
symplectic duality \cite{BLPW16}, Hikita formulated a fixed-scheme
realization of cohomology and established it for the Hilbert scheme of
points in the plane, with related type-$A$ Spaltenstein and smooth
hypertoric examples \cite{Hikita17}.  These examples explain why the
infinitesimal structure of the fixed locus, rather than its set of
points, is central to the question.

Nakajima proposed a mathematical Coulomb-branch construction for
three-dimensional gauge theories \cite{Nakajima16};
Braverman--Finkelberg--Nakajima constructed its commutative convolution
algebra \cite{BFN18,BFN19}, extending the affine-Grassmannian convolution
framework of Bezrukavnikov--Finkelberg--Mirkovi\'c \cite{BFM05}.
For quiver gauge theories the corresponding
Higgs quotient is a Nakajima variety \cite{Nakajima94}.
The coefficient-compatible formulation has important predecessors.
Kamnitzer--Tingley--Webster--Weekes--Yacobi record Nakajima's equivariant
conjecture in the affine-Grassmannian setting and strengthen its
specialized form by requiring compatibility with tautological classes
\cite[Conjectures~8.9--8.10]{KTWWY19}.
Hilburn--Kamnitzer--Weekes formulate the equivariant quantized refinement
over the common coefficient algebra
\cite[Section~5.7, Conjecture~5.26]{HKW23}.
Krylov--Shlykov state the general-quiver comparison for the
\emph{full topological-torus fixed scheme} as equality of coefficient-map
kernels \cite[Section~8.1, Conjecture~8.1]{KrylovShlykov25}.
We first establish that equality and then compare the full-torus fixed
scheme with the stability-cocharacter fixed scheme.
Thus compatibility with the coefficient maps is part of the conjectural
formulation itself.

Krylov--Shlykov prove the Gieseker, or framed Jordan-quiver, case
\cite{KrylovShlykov25}.  Dumanski--Krylov give the formulation used here
and compare cohomological, equivariant, and $K$-theoretic versions
\cite{DK25}.  Their finite-type conclusion includes types $A,D,E_6$;
in types $E_7,E_8$ it imposes the framing restrictions in
\cite[Corollary~2.17]{DK25}.  Their Remark~2.18 explains why the earlier
unrestricted ADE assertion must be qualified.
Recent developments also include Setiabrata's Hikita surjectivity result
for the base-affine-cotangent example \cite{Setiabrata24}, and the refined quantum
Hikita theorem of Dinkins--Karpov--Krylov for ADE quivers with
minuscule-supported framing and for the Jordan quiver
\cite[Theorem~B]{DinkinsKarpovKrylov26}.
Theorem~\ref{thm:main} treats the cohomological algebra for all the finite
quiver data above.  It does not assert the quantized statement in
\cite[Equation~(2.2)]{DK25}, nor extend the theorem to arbitrary
reductive gauge groups; outside this setting, proposed Hikita-type
identifications can fail \cite[Example~6.7]{HoangKrylovMatvieievskyi24}.

Several established tools enter the argument in distinct ways.
BFN abelianization turns minuscule convolution classes into explicit
Laurent expressions, and Weekes's generation theorem makes these classes
sufficient for arbitrary quivers \cite{BFN18,BFN19,Weekes19}.
McGerty--Nevins's ordinary and equivariant Kirwan surjectivity results
supply the tautological presentation on the Higgs side
\cite[Theorem~1.2 and Corollary~1.5]{MN18}.
Hausel's arithmetic Fourier-transform formula computes the Betti numbers
that eventually determine the size of the coefficient quotient
\cite{Hausel06}.  Character-linearized GIT and the topological
Kempf--Ness correspondence connect regular stability to the flat
hyperk\"ahler quotient construction
\cite{King94,Hoskins14,HKLR87,Nakajima94}.
Our use of these inputs is explained at their precise interfaces below.

The finite-dimensional critical-model viewpoint is related to Nakajima's
original construction and to Botta--Tamagni's critical lattice and
quasimap models \cite{Nakajima16,BottaTamagni26}.  Here the specific
stable-pair test proves an inclusion of coefficient ideals; it does not
require their entire Coulomb-action construction.  Deformation arguments
also appear in Krylov--Shlykov's Gieseker proof
\cite[Remark~1.15]{KrylovShlykov25}.  Our loop-removal step uses subgroup
localization \cite{Brion98}, with the prequotient perspective of
\cite[Proposition~2.3]{DK25}, and an explicit descent for homogeneous
ideals.  These connections distinguish the established methods from the
particular comparison arguments carried out here.

The difficulty is not just to match dimensions.  One must identify the
relations under the specified maps, with their grading and nilpotents,
for arbitrary flavors and for every regular stability.  Neither a
calculation on closed fixed points nor an isomorphism after inverting
generic parameters establishes that assertion.  A further distinction is
that a regular character can satisfy $\nu\cdot\mathbf b=0$ for an integer
vector $0<\mathbf b\leq\mathbf v$.  Therefore the full topological-torus calculation
does not by itself identify the cocharacter-fixed quotient.

\subsection{Proof strategy and organization}

Define the intermediate coefficient algebra and the two ideals to be
compared by
\begin{equation}\label{eq:intro-two-ideals}
 \begin{split}
 B_F&=\C[Y_{\mathfrak f}^{T_{\mathrm{top}}}]
      =A_F/\langle A_F^{(d)}:d\ne0\rangle,\\
 I_{\mathrm{top}}&=\ker(R\longrightarrow B_F),\qquad
 I_H=\ker(R\xrightarrow{\phi_H}H_F^*(X)),
 \end{split}
\end{equation}
where $A_F^{(d)}$ is the subspace of topological charge
$d\in\pi_1(G)$; its cocharacter weight is $\nu\cdot d$.  Neither of
these is the cohomological degree used in Hilbert series.
Section~\ref{sec:preliminaries} explains why the
coefficient maps to $B_F$ and to $H_F^*(X)$ are surjective.
Thus the first task is precisely the equality
$I_{\mathrm{top}}=I_H$.  The second task is the equality of the
cocharacter and full-torus defining ideals inside $A_F$.
A subscript $0$ means that the flavor torus is trivial, so
$R_0=H_G^*(\mathrm{pt})$ and $B_0$ denotes the corresponding full-torus
quotient.  Figure~\ref{fig:proof-route} separates the two tasks.

\paragraph{1. Stable pairs give a quotient map.}
Sections~\ref{sec:correspondences} and \ref{sec:stable-pairs} prove
$I_{\mathrm{top}}\subseteq I_H$ for all the data of the theorem.
The nontrivial minuscule generators used here are dressed lattice
modifications whose coweight entries are all in $\{0,1\}$ or all in
$\{0,-1\}$, with nonzero topological charge; a subset may involve several
vertices.  A convolution word is a finite product of these classes.
By minuscule generation, a coefficient in $I_{\mathrm{top}}$ is a finite
sum of positive-length convolution words, each of total topological charge zero.
Section~\ref{sec:correspondences} realizes each such word by
finite-dimensional lattice correspondences, keeping all degree shifts.
The test in Section~\ref{sec:stable-pairs} evaluates them on stable
cotangent pairs.  The stability character assigns an additive degree to
each modification, so the degrees in a neutral word sum to zero.
After reordering each word, its first modification is nontrivial and has
nonpositive character-line degree; stability excludes it.  On the identity
correspondence, the same test gives a split restriction onto $H_F^*(X)$,
which detects the coefficient class itself.  The tested relation therefore vanishes
in cohomology, yielding the graded coefficient-compatible surjection
$B_F\twoheadrightarrow H_F^*(X)$ of
Proposition~\ref{prop:kernel-inclusion}.

\paragraph{2. Added loops make the size comparison explicit.}
Temporarily suppose that the quiver has a loop at every vertex of
positive gauge dimension, take no flavor, and use determinant stability
$\nu_{\det}=\prod_i\det_{V_i}$.
The loop weights cancel the Grassmannian denominators in minuscule
abelianization.  Section~\ref{sec:laurent} proves a Laurent model with
polynomial coefficients and filters $B_0$ by multipartitions
of the gauge dimension vector.  Its successive quotients are bounded by
invariant parts of polynomial quotients defined by cut products: products of framing variables
and differences across a subset of slots and its complement.
Section~\ref{sec:counting} computes their Hilbert series
using a reachability stratification and identifies the resulting sum
with Hausel's ordinary Poincar\'e polynomial.  The upper bound and the
surjection from the first step force equality in every degree.  This is
an isomorphism under the original coefficient map, not merely a numerical
identity; see Corollary~\ref{cor:loop-isomorphism}.

\paragraph{3. Deformation returns to the original data.}
For an arbitrary quiver $Q$, let $Q^+$ be obtained by adding one loop at
each vertex of positive gauge dimension.  Give the new loops the common
weight one of an auxiliary flavor torus $H=\C^\times$, with parameter
$z\in H_H^2(\mathrm{pt})$.
The comparison in Section~\ref{sec:deformations} uses a lift of the
no-flavor isomorphism for $Q^+$ to this flavor.  Near zero gauge parameters and nonzero mass
$z=1$, the new minuscule Euler factors are invertible on the Coulomb
side.  Subgroup localization on the stable prequotient gives the
corresponding comparison on the Higgs side.  The resulting equality of
localized coefficient ideals descends to the original homogeneous ideals
in $R_0$: their equality at the graded origin determines them globally.
This homogeneity argument, not specialization of an unspecified generic
isomorphism, is the content of Proposition~\ref{prop:loop-removal}.
Local constancy of regular stability, quaternionic rotation, and generic
constancy of compactly supported Betti numbers then compare all regular
characters.  A final flavor lift proves
$B_F\simeq H_F^*(X)$ in Theorem~\ref{thm:full-torus}.

\paragraph{4. Complementary monopoles identify the cocharacter quotient.}
Section~\ref{sec:cocharacters} proves, independently of the preceding
cohomology comparison, that the natural map
$\C[Y_{\mathfrak f}^{\nu}]\to B_F$ is an isomorphism.
It combines polynomial-form Hall-algebra generation
\cite{JindalNegut26} with the monopole maps of
Lemma~\ref{lem:polynomial-monopole-maps}, derived from
\cite[Theorem~5.7]{Tsymbaliuk23Difference} with explicit product signs.
The resulting finite identities can be specialized without dividing by
rotation or arrow parameters.
After adjoining a rank-one framing vertex, any surviving nontrivial minuscule
monopole has a surviving complementary monopole.  Hall generation would
then decompose both charges into indecomposable charges of the enlarged
quiver, all of weight zero for the extended character.
Lemma~\ref{lem:root-moment} realizes these charges by moment-level blocks.
Their semistable direct sum has a stabilizer that remains
positive-dimensional even after removing the ineffective diagonal
scalar.  This contradicts the freeness in regularity, including when
the semistable locus is empty.  Proposition~\ref{prop:cocharacter-fixed}
therefore identifies the two defining ideals, completing the
scheme-theoretic comparison.

\begin{figure}[hbp]
\centering
\begingroup
\small
\setlength{\fboxsep}{7pt}
\begin{minipage}[t]{0.47\textwidth}
\centering
\textbf{The coefficient-ideal comparison}\par\medskip
\fbox{\parbox{0.89\linewidth}{\centering
Every $Q,F$ and regular $\nu$\\
Stable-pair test\\
$I_{\mathrm{top}}\subseteq I_H$\\
Sections~\ref{sec:correspondences}--\ref{sec:stable-pairs}}}
\par\smallskip$\Downarrow$\par\smallskip
\fbox{\parbox{0.89\linewidth}{\centering
$Q^+$, no flavor, determinant stability\\
Laurent model and cut-ideal count\\
$B_0(Q^+)\simeq H^*(X_{\det}(Q^+))$\\
Sections~\ref{sec:laurent}--\ref{sec:counting}}}
\par\smallskip$\Downarrow$\par\smallskip
\fbox{\parbox{0.89\linewidth}{\centering
Lift to the added-loop flavor $H$\\
Compare near gauge $0$, mass $z=1$\\
Descend homogeneous ideals for $Q$\\
Proposition~\ref{prop:loop-removal}}}
\par\smallskip$\Downarrow$\par\smallskip
\fbox{\parbox{0.89\linewidth}{\centering
All regular stabilities, then all flavors\\
$B_F\simeq H_F^*(X)$\\
Theorem~\ref{thm:full-torus}}}
\end{minipage}
\hfill
\begin{minipage}[t]{0.47\textwidth}
\centering
\textbf{The fixed-scheme comparison}\par\medskip
\fbox{\parbox{0.89\linewidth}{\centering
Suppose a nontrivial minuscule class\\
survives in the cocharacter quotient.\\
Adjoin a rank-one framing vertex;\\
its complementary class also survives}}
\par\smallskip$\Downarrow$\par\smallskip
\fbox{\parbox{0.89\linewidth}{\centering
Polynomial Hall generation\\
Both surviving charges decompose\\
into indecomposable charges of the\\
enlarged quiver, all of weight zero\\
Lemma~\ref{lem:root-charge-generation}}}
\par\smallskip$\Downarrow$\par\smallskip
\fbox{\parbox{0.89\linewidth}{\centering
Lemma~\ref{lem:root-moment} gives moment-level blocks.\\
Their direct sum retains a\\
positive-dimensional stabilizer after\\
removing the diagonal scalar,\\
contradicting regularity}}
\par\smallskip$\Downarrow$\par\smallskip
\fbox{\parbox{0.89\linewidth}{\centering
$\C[Y_{\mathfrak f}^{\nu}]\simeq B_F$\\
as quotients of $A_F$\\
Proposition~\ref{prop:cocharacter-fixed}}}
\end{minipage}
\par\bigskip
$\displaystyle
 \C[Y_{\mathfrak f}^{\nu}]\ \xrightarrow{\ \sim\ }\
 B_F\ \xrightarrow{\ \sim\ }\ H_F^*(X)$
\endgroup
\caption{The two comparisons in the proof.  Here $Q^+$ adds one loop at
each positive-rank vertex of $Q$; dimension and framing vectors are
unchanged and suppressed.  The left column uses the stable-pair
surjection in the dimension comparisons.  The right column is
independent of the left and identifies defining ideals, not just closed
fixed points.  Both final maps preserve the coefficient ring and grading.}
\label{fig:proof-route}
\end{figure}

